\subsection{Two Reflective Towers, Two Kinds of Control}

The Hybrid Group Architecture distinguishes reflection on an individual
object from reflection on the group in which objects execute.
Masuhara's thesis, Section 3.3, develops this distinction~\cite{masuhara}.

\begin{longtable}{L{0.26\textwidth}L{0.32\textwidth}L{0.32\textwidth}}
\textbf{Level} & \textbf{What it controls} & \textbf{Example}\\\hline
Individual metaobject & The object's state environment, scripts, and
normal message queue. & Inspect a counter's state or change its evaluator.\\
Group kernel objects & Shared evaluation and object creation policy. &
Give a group a custom evaluator for a new control form.\\
Higher meta-levels & The corresponding behavior of metaobjects and
meta-groups. & Reflect on the evaluator itself.\\
\end{longtable}

This conceptual hierarchy is not allocated as an infinite collection
of objects. Chapter 4 of the thesis explains lazy creation and lightweight
metaobjects: forwarding an ordinary message should not force expensive
reification. Section 4.7 discusses non-reifying objects; Chapter 7
summarizes the use of partial compilation~\cite{masuhara}.
Sequential host computations can be compiled while selected operations
still pass through the reflective protocols.

\subsection{Inspecting an Object's State}

Load the reflection example in a fresh session:

\begin{lstlisting}
(inc "tutorial/examples/05-reflection" :use-source t)
[observed-counter <== :inc]                  ; => 1
[[[meta observed-counter] <== :state]
 <== [:value-of 'n]]                         ; => 1
[[[meta observed-counter] <== :state]
 <== [:set 'n 40]]                           ; => 40
[observed-counter <== :inc]                  ; => 41
\end{lstlisting}

The steps are significant. \code{[meta observed-counter]} obtains the
counter's metaobject. The \code{:state} query returns an environment
object, not a plain property map. Sending \code{[:value-of 'n]} to that
environment retrieves a binding; \code{[:set 'n 40]} changes it.
The quoted symbol is the variable's name as data.
The next ordinary increment observes that change: the meta-level
representation is connected to the computation it describes.

This bypasses the object's application protocol. Use it to understand
reflection and debugging; ordinary clients should continue to use
the counter's defined messages. Default metaobjects also answer
\code{:scripts}, \code{:queue}, and \code{:evaluator}.
Their precise replies are runtime objects or internal queue entries.
The operation \code{[den metaobject]} returns the associated lower-level
object, called its \emph{denotation}. It is not a copy or a variable's
value. In the original inheritance project we will check both identity
directions and distinguish this relationship from group membership.

\subsection{Defining a Lazy Control Form}

An ordinary function evaluates its arguments before receiving them.
A control form must instead decide which expression to evaluate.
The next example implements a choice using the group's evaluator:

\wholelisting[linerange={9-18,25-26}]{tutorial/examples/05-reflection.abclr2}

The new evaluator handles the \code{:choose} protocol and forwards all
other messages to a primary evaluator. The internal helper
\code{r2sys::create-primary-eval} matches the helper used by the
repository's examples; it is specific to this implementation.
Inside a group definition, \code{group} names the group being constructed.

The expression variables \code{predicate}, \code{then}, and \code{else}
contain compiled expression representations. They are not ordinary
argument results. The \code{[:do expression Env Id Gid Eval]} protocol
asks the evaluator to execute one representation in its original context.
Read the context fields by role, not as application arguments:

\noindent\begin{minipage}{\linewidth}
\begin{longtable}{L{0.13\textwidth}L{0.77\textwidth}}
\textbf{Field} & \textbf{Meaning in this evaluation request}\\\hline
\code{Env} & Environment used for variable lookup.\\
\code{Id} & Metaobject identity of the object whose expression is being
executed; evaluating \code{Me} produces \code{[den Id]}.\\
\code{Gid} & Group manager for that execution context; evaluating
\code{group} produces this value.\\
\code{Eval} & Evaluator used for the expression's evaluation operations.\\
\code{C} & Continuation receiving this expression's result, not
necessarily the application's original caller.\\
\end{longtable}
\end{minipage}

When a metaobject's own handler executes, \code{Me} is that metaobject,
not the application object it manages. Likewise, \code{Gid} is not
automatically the group of \code{Id}: a metaobject belongs to a group
at its own level. Keep object, metaobject, and group relationships separate.

The first send evaluates the predicate. Its continuation selects
one branch representation and forwards the final destination \code{C}.
No unselected branch is evaluated.

\sourcelisting[linerange={28-36,41-41}]{tutorial/examples/05-reflection.abclr2}

\begin{codeblock}
\begin{lstlisting}
[choice-test <== [:choose 3]]   ; => 2
[choice-test <== :selected]     ; => :ODD
[choice-test <== [:choose 8]]   ; => 4
[choice-test <== :selected]     ; => :EVEN
\end{lstlisting}
\end{codeblock}

The current evaluator receives compiled argument expressions through
\code{reify-args}. The corresponding form for the current object's
metaobject is \code{reify-args-to-meta}. Both preserve the evaluation context,
and neither is available in a non-reifying object's script.
This example uses a reflective control form to add behavior; it does
not change the meaning of the ordinary \code{if} operator.

The tag supplied to \code{reify-args} is evaluated first. The remaining
argument expressions are passed as compiled representations in a message
of the form \code{[Tag Expressions Env Id Gid Eval] @ C}; they have
not already been reduced to argument values~\cite{guide}.

\paragraph{Exercise 5.}
Extend the choice group with a \code{:unless} protocol that evaluates
one body only when its predicate is false. Forward the final result
to the original continuation in both cases.
Check both branches against Section~\ref{sol:unless}.

\subsection{Extension Parameters and Host Events}

The user's guide identifies additional parameters on message transmission,
object creation, and assignment as extension points~\cite{guide}.
They carry information to the execution machinery; they are not extra
elements of the application's message payload. For example, in the
schematic send \code{[target <= payload @ reply :RVT stamp]}, the
evaluator representation is
\code{[:send-past Target Payload Reply :RVT Stamp]}.
Capitalized names here denote evaluated values.
It is the selected evaluator's responsibility to interpret this extension.
An ordinary receiver still receives \code{Payload}.

Likewise, \code{(priority n)} in an object definition becomes a
\code{:priority} argument in its creation representation, rather than
automatically creating a state variable or a universal scheduling rule.
The priority-control examples supply machinery that gives this parameter
meaning. An extended assignment \code{[x := value tag argument]}
passes the evaluated extras to the environment's \code{:set} request.
Custom environment objects can therefore participate in update policy.
Do not assume the default environment acts on an arbitrary tag.

Two library facilities illustrate this principle without extending the
core syntax. Script-pattern \code{where} constraints require
\code{heavy-meta}; constraints inside an ordinary \code{match} form
remain unsupported. Multicast requires \code{multicast} and a suitable
evaluator created by \code{multicast-eval-gen}; merely loading the
library does not replace every object's evaluator.
The original \filepath{lib/multicast.abclr2} constructs a template for
nested target lists and collects replies into their corresponding
positions. This is a library-defined protocol, not the unsupported
brace-delimited synchronous-send form.
Section~\ref{sec:multicast} works through its original example and reply paths.
These libraries need their own setup and reply/completion analysis,
just as the express library does.

The guide also describes a boundary for host events~\cite{guide}.
At Lisp level, \code{r2sys:send-signal} delivers its arguments as a
\code{[:signal ...]} message to the object bound to
\code{*signal-handler*}. The default handler ignores events; a
\code{nil} handler disables delivery.
Alternatively, the host can set \code{r2sys:*signal-request*};
the runtime checks and clears it before forwarding the request.
This is a single pending-request slot, not a general event queue.
The ABCL side handles an event as a message. It does not follow that
every operating-system signal or timer is automatically connected to
this interface: that requires host-specific integration.
The interface also does not promise express preemption.

\subsection{A Safe Experiment with Synthetic Host Events}
\label{sec:host-events}

The event interface turns a host notification into an ABCL message.
You can explore that boundary without a timer, an operating-system signal,
or another thread. The following supplemental exercise uses the interface
described by the user's guide and implemented in \filepath{src/signal.lsp};
there is no corresponding runnable event example in \filepath{examples/}.
It changes no original source and installs no operating-system callbacks.

\paragraph{Receive and inspect an event.}
The file \filepath{tutorial/examples/14-host-events.abclr2} defines an
ordinary recording object. Loading it does not install it as a handler:

\wholelisting[firstline=2,lastline=8]{tutorial/examples/14-host-events.abclr2}

The dotted pattern \code{[:signal . parameters]} binds the argument tail
as a list. Each saved record contains that list, the reply destination,
and the sender. A host notification arrives with both reply and sender
\code{nil}: it is not an application call waiting for \code{!}.
The separate \code{:events} query does reply, allowing us to inspect
what the handler has processed.

Start a fresh interpreter at the repository root and save its current
handler before temporarily replacing it:

\begin{codeblock}
\begin{lstlisting}
(inc "tutorial/examples/14-host-events" :use-source t)
[saved-event-handler := *signal-handler*]
[*signal-handler* := event-inbox]
(r2sys:send-signal :tick 7)
[event-inbox <== :events]           ; => (((:TICK 7) NIL NIL))
\end{lstlisting}
\end{codeblock}

The \code{r2sys:} prefix names the exported host interface in the
ABCL/R2 system package. \code{send-signal} accepts an ordinary Lisp
\code{\&rest} argument list and constructs \code{(:signal :tick 7)}.
It submits the notification; returning from that function is not an
acknowledgment from your handler. This finite interactive exercise lets
the runtime dispatch the event before inspecting the inbox. Do not infer
a general ordering guarantee between independent producers, or use a
notification as proof that unrelated work finished.

Passing one list is different from passing its elements as arguments:

\begin{codeblock}
\begin{lstlisting}
[event-inbox <== :clear]
(r2sys:send-signal '(:tick 7))
[event-inbox <== :events]           ; => ((((:TICK 7)) NIL NIL))
\end{lstlisting}
\end{codeblock}

The parameter list now has one element, itself a list. Similarly,
\code{(r2sys:send-signal)} sends an event with an empty parameter list,
while \code{(r2sys:send-signal nil)} sends one with parameter list
\code{(nil)}. Neither call means ``disable events.''

\paragraph{Explore the single pending-request slot.}
The other interface is a \emph{host Lisp variable},
\code{r2sys:*signal-request*}. It is distinct from a variable stored in
an ABCL environment. An ABCL assignment such as
\code{[r2sys:*signal-request* := value]} would not be the host-side
write this interface requires. The teaching file supplies small host
functions to make that boundary explicit:

\wholelisting[firstline=11,lastline=17]{tutorial/examples/14-host-events.abclr2}

These function bodies run in host Lisp. Their \code{setf} writes the
actual host variable; \code{demo-event-slot} reads it. They are teaching
helpers, not additional built-in ABCL/R2 operations.
In the same fresh session, first check that the slot is empty:

\begin{codeblock}
\begin{lstlisting}
(demo-event-slot)                  ; => NIL
[event-inbox <== :clear]
(demo-post-event '(:tick 7))
[event-inbox <== :events]           ; => ((((:TICK 7)) NIL NIL))
(demo-event-slot)                  ; => NIL: consumed by the runtime
\end{lstlisting}
\end{codeblock}

A non-nil slot value becomes \emph{one} argument following \code{:signal}.
A list stored there is not spliced into multiple arguments.
The original polling function makes this explicit:

\wholelisting[firstline=57,lastline=61]{src/signal.lsp}

It saves the value, clears the slot, and forwards \code{(:signal value)}.
The scheduler calls this function at the top of its dispatch loop in
\filepath{src/process.lsp}. The slot is not a background service that
keeps an otherwise idle interpreter polling, and its check does not
interrupt a host Lisp computation halfway through execution.

Two writes before that polling point do not create two pending events:

\begin{codeblock}
\begin{lstlisting}
[event-inbox <== :clear]
(demo-replace-event :first :second)
[event-inbox <== :events]           ; => (((:SECOND) NIL NIL))
\end{lstlisting}
\end{codeblock}

Both writes happen inside one host function, with no ABCL scheduler
step between them. Only the last value remains to be polled. This
demonstrates overwrite, not queueing or reliable event aggregation.
Writing \code{nil} means no pending request; use
\code{(r2sys:send-signal nil)} if you actually need an event whose
argument is \code{nil}. These helpers assume an isolated fresh session;
do not overwrite a pending request owned by another event producer.

\paragraph{Disable delivery, then restore the saved handler.}
The handler binding controls both interfaces:

\begin{codeblock}
\begin{lstlisting}
[event-inbox <== :clear]
[*signal-handler* := nil]
(r2sys:send-signal :dropped)
(demo-post-event :also-dropped)
(demo-event-slot)                  ; => NIL after polling
[*signal-handler* := event-inbox]
[event-inbox <== :events]           ; => NIL: no replay
[*signal-handler* := saved-event-handler]
(eq *signal-handler* saved-event-handler) ; => T
\end{lstlisting}
\end{codeblock}

Disabling delivery does not buffer notifications for later. A polled
slot value is cleared even when no handler is installed. Conversely,
changing the handler binding is not a cancellation of messages already
sent to a previous handler. Always restore the object you saved, rather
than assuming that the default handler is still the desired one.
The normal default handler prints an ignored-signal diagnostic; ``ignored''
does not necessarily mean silent.

No step in this exercise enables keyboard interrupts or timer delivery.
The historical keyboard-interrupt wrapper in \filepath{src/signal.lsp}
contains host-specific branches; its fallback simply executes its body.
Do not infer an SBCL operating-system hook from the presence of this API.
Nor is this an Express-message experiment: the original Interrupt and
Express examples have their own reflective execution protocols.

\paragraph{Evidence and exercise.}
Run \code{make -C tutorial host-events-check}. Twenty-five probes in two
fresh processes check handler preservation, direct argument shapes,
nil sender/reply fields, slot consumption and overwrite, disabled delivery,
absence of replay, and restoration. Inputs, copied teaching source and
transcripts remain in \filepath{build/host-event-tests/}. Tests install
no operating-system callbacks and modify no runtime or original sources.

Why does a list-valued slot notification differ from
\code{(r2sys:send-signal :tick 7)}? \emph{Answer:} The slot carries one
value; the function accepts several arguments. Why is a zero-valued
slot request still delivered? \emph{Answer:} Only \code{nil} is false
in Lisp; zero is non-nil. Why should the signal handler not rely on
\code{!} to acknowledge the host? \emph{Answer:} The notification has
no application reply destination. Define a separate acknowledgment
protocol if your application needs one.

